\documentclass[a4paper,12pt]{article}
\usepackage[utf8]{inputenc}
\usepackage[italian]{babel}
\usepackage{listings}
\usepackage{xcolor}
\usepackage{float}
\usepackage{placeins}
\usepackage{url}
\setlength{\emergencystretch}{3em}

\lstset{
  language=C,
  backgroundcolor=\color{lightgray!20},
  frame=single,
  breaklines=true,
  numbers=left,
  numberstyle=\tiny,
  stepnumber=1,
  numbersep=5pt,
  captionpos=b,
  basicstyle=\ttfamily\small,
  keywordstyle=\color{blue}\bfseries,
  commentstyle=\color{green!60!black},
  stringstyle=\color{red},
  showstringspaces=false
}

\title{LabReSiD26 -- Hands-On 6}
\author{Davide Ferrara -- 518629}
\date{}

\begin{document}

\maketitle

% ---------------------------------------------------------------------------
\section{Esercizio 1: Analisi di \texttt{server\_epoll.c}}
% ---------------------------------------------------------------------------

\subsection*{Domanda 1 -- Differenza strutturale tra \texttt{server\_epoll.c} e \texttt{server\_select\_demo.c}}

La differenza strutturale più evidente è l'eliminazione della rubrica \texttt{client\_fds[]},
di \texttt{max\_fd} e della ricostruzione del set ad ogni ciclo con
\texttt{FD\_ZERO}/\texttt{FD\_SET}. In \texttt{server\_select\_demo.c} il programma
deve mantenere queste strutture manualmente; in \texttt{server\_epoll\_demo.c} è il
kernel a tracciare i fd registrati tramite il suo albero rosso-nero interno.
\texttt{max\_fd} non è più necessario perché \texttt{epoll\_wait()} non scansiona un
range di fd: restituisce direttamente i soli fd pronti nell'array
\texttt{events[]}.

\medskip
\noindent\textbf{Prima (select):}
\begin{lstlisting}
/* rubrica da mantenere manualmente */
int client_fds[MAX_CLIENTS];
for (int i = 0; i < MAX_CLIENTS; i++) client_fds[i] = -1;
int max_fd = server_fd;

/* ad ogni iterazione si ricostruisce il set */
fd_set read_set;
FD_ZERO(&read_set);
FD_SET(server_fd, &read_set);
for (int i = 0; i < MAX_CLIENTS; i++)
    if (client_fds[i] != -1) FD_SET(client_fds[i], &read_set);
int n = select(max_fd + 1, &read_set, NULL, NULL, NULL);
\end{lstlisting}

\noindent\textbf{Dopo (epoll):}
\begin{lstlisting}
/* nessuna rubrica, nessun max_fd */
int epfd = epoll_create1(0);
struct epoll_event events[MAX_EVENTS];

/* ad ogni iterazione: epoll_wait restituisce solo i fd pronti */
int n = epoll_wait(epfd, events, MAX_EVENTS, -1);
for (int i = 0; i < n; i++) {
    int fd = events[i].data.fd;
}
\end{lstlisting}

\subsection*{Domanda 2 -- Dove sono finiti \texttt{client\_fds[]} e \texttt{max\_fd}}

Entrambe le strutture spariscono dal codice utente. Se ne occupa il kernel: ogni
\texttt{epoll\_ctl(EPOLL\_CTL\_ADD, fd, ...)} inserisce il fd nell'albero rosso-nero
interno dell'istanza epoll.

\texttt{epoll\_wait()} si limita a copiare i fd pronti nell'array \texttt{events[]}
 con costo O(eventi pronti).

\texttt{max\_fd} non esiste più perché \texttt{epoll\_wait()} non riceve un range di fd
da scansionare ma conosce già i fd registrati tramite l'albero rosso-nero e consegna
direttamente ciò che è pronto.

\subsection*{Domanda 3 -- \texttt{epoll\_wait()} ritorna con \texttt{N > 1}}

Dal manuale:

\begin{verbatim}
  On success, epoll_wait() returns the number of file
  descriptors ready for the requested I/O operation,
  or zero if no file descriptor became ready
  during the requested timeout milliseconds.
\end{verbatim}

\texttt{epoll\_wait()} può quindi ritornare con \texttt{n > 1} quando più fd diventano
pronti prima che il server torni a chiamarlo. Nel caso dello script fornito
per aumentare la probabilità che due o più client si sovrappongano è necessario
diminuire il delay tramite il secondo parametro dello script:

\begin{lstlisting}
./test_multifd.sh 10 0
\end{lstlisting}

Con delay=0 le connessioni avvengono immediatamente e con un numero alto di client
aumenta la probabilità che più eventi si accodino nella ready list prima della prossima
chiamata a \texttt{epoll\_wait()}, ottenendo \texttt{n > 1}.

\subsection*{Domanda 4 -- \texttt{epoll\_ctl(DEL)} prima o dopo \texttt{close(fd)}}

\texttt{epoll\_ctl(DEL)} va chiamato \textbf{prima} di \texttt{close(fd)}. Dal kernel
2.6.9 la rimozione esplicita è diventata opzionale (il kernel rimuove automaticamente
il fd alla chiusura), ma l'ordine corretto resta DEL poi close per due motivi.

\textbf{Gestione degli errori.} Dal manuale: \textit{``When successful, epoll\_ctl()
returns zero. When an error occurs, epoll\_ctl() returns -1 and errno is set to
indicate the error.''} Se si chiude prima il fd, questo non è più valido: una
successiva \texttt{epoll\_ctl(DEL)} fallirebbe con \texttt{EBADF} (Bad File
Descriptr) e non sarebbe possibile gestire l'errore.

\textbf{Rischio di fd reuse.} Alla \texttt{close()}, il kernel libera immediatamente
il numero di fd. Un altro thread o la stessa \texttt{accept()} potrebbe ricevere
quel numero come nuovo fd. Chiamare \texttt{epoll\_ctl(DEL, fd)} dopo rimuoverebbe
il fd \emph{nuovo} dall'istanza epoll, non quello vecchio, causando un comportamento
indefinito silenzioso senza alcun errore visibile.

L'ordine corretto è quindi:
\begin{lstlisting}
epoll_ctl(epfd, EPOLL_CTL_DEL, fd, NULL);
close(fd);
\end{lstlisting}

\subsection*{Domanda 5 -- \texttt{epoll\_create1()} chiamato una volta sola}

\texttt{epoll\_create1()} crea un'istanza epoll nel kernel e restituisce un file
descriptor che la rappresenta (\texttt{epfd}). Chiamarlo ad ogni iterazione del loop
causa tre problemi.

\textbf{File descriptor leak.} Ogni chiamata apre un nuovo \texttt{epfd} senza che
il precedente venga mai chiuso. Il kernel mantiene allocati albero rosso-nero e ready
list per ogni istanza: iterando abbastanza a lungo si esauriscono i fd disponibili
del processo (solitamente 1024 massimi per processo).

\begin{lstlisting}
ulimit -n
1024 /* Numero massimo di FD nel mio sistema*/
\end{lstlisting}

\textbf{Memoria kernel non liberata.} Le strutture interne di ogni istanza epoll
(albero rosso-nero + ready list) rimangono allocate fino alla \texttt{close(epfd)}.
Con un nuovo \texttt{epfd} ad ogni ciclo e nessuna chiusura, la memoria kernel cresce
senza limite causando possibili OOM.

\textbf{Perdita delle registrazioni.} Ogni nuova istanza parte vuota: tutti i client
già connessi e registrati con \texttt{epoll\_ctl(EPOLL\_CTL\_ADD)} spariscono. Il server
perderebbe traccia di tutte le connessioni attive ad ogni iterazione.

\subsection*{Domanda 6 -- Confronto riga di compilazione: \texttt{epoll} vs \texttt{io\_uring}}

Un programma che usa \texttt{io\_uring} tramite l'API ad alto livello richiede
\textbf{liburing} (\texttt{-luring}), la libreria userspace di Jens Axboe.

Un esempio reale verificato è \textbf{QEMU}: il suo \texttt{meson.build}\footnote{\url{https://github.com/qemu/qemu/blob/master/meson.build}} dichiara
esplicitamente la dipendenza da liburing tramite pkg-config:

\begin{verbatim}[meson.build]
linux_io_uring = dependency('liburing', version: '>=0.3',
                            required: get_option('linux_io_uring'),
                            method: 'pkg-config')
\end{verbatim}

Questo produce in compilazione il flag \texttt{-luring}.

\medskip
\noindent Confronto righe di compilazione:

\begin{lstlisting}[language=bash]
# server_epoll.c con nessuna libreria esterna
gcc -Wall -O2 -o server_epoll server_epoll.c

# programma con link di uring
gcc -Wall -O2 -o server_uring server_uring.c -luring
\end{lstlisting}

\end{document}
